\subsection{Why RouterBSI Trails SkillRouter on Two Datasets}
\label{app:routing-gap}
The mean per-retriever relative nDCG@10 shortfalls against SkillRouter are 5.01\% on MedCalc-Bench and 6.10\% on TheoremQA, using SkillRouter as the denominator. To investigate these gaps, we match the same 165 MedCalc-Bench and 105 TheoremQA test tasks across BM25, Vector, and SkillSight. Both methods receive identical ordered top-50 candidate lists in all 810 task--retriever observations. A \emph{ranking-loss case} is an observation where SkillRouter obtains higher nDCG@10, not a downstream answer failure. Pooling retrievers repeats questions rather than providing independent training replications.

\begin{table*}[t]
\centering
\small
\begin{tabular}{lrr}
\toprule
Quantity & MedCalc-Bench & TheoremQA \\
\midrule
Task--retriever observations & 495 & 315 \\
SkillRouter-better ranking cases & 35 & 66 \\
Distinct tasks among these cases & 28 & 46 \\
RouterBSI-better ranking cases & 6 & 32 \\
Tied ranking scores & 454 & 217 \\
Loss cases: required skill not retrieved & 0 & 0 \\
Loss cases: required skill retrieved but rejected & 35 (100\%) & 55 (83.33\%) \\
Loss cases: required skill selected but ranked lower & 0 & 11 \\
Gross ranking loss in retrieved-but-rejected cases & 100\% & 86.38\% \\
\bottomrule
\end{tabular}
\caption{Matched nDCG@10 comparisons. Loss-case percentages use the 35 and 66 SkillRouter-better observations, respectively. Gross ranking loss sums positive per-observation SkillRouter-minus-RouterBSI score differences before offsetting RouterBSI's gains; its shares are not shares of the net gap or estimated recoverable improvements. These are retrospective diagnostics, not intervention results.}
\label{tab:app-routing-gap}
\end{table*}

\paragraph{Required-skill nonactivation.}
Table~\ref{tab:app-routing-gap} identifies required-skill rejection as the dominant observed bottleneck. All 35 MedCalc-Bench ranking-loss cases have an empty selected set and all candidate predicates false. The final ranking therefore remains the raw retrieval order, while SkillRouter places the required calculator first: RouterBSI fails to promote an available skill rather than demoting it below selected competitors. Across all MedCalc-Bench observations, 76 retrieved required skills are rejected, but 41 remain at rank one and incur no comparative ranking loss. In TheoremQA, 55 of 66 loss cases reject the required skill; 42 of these place selected nonrequired skills above it. The other 11 loss cases select the required skill but rank it below another selected skill. Retrieval misses do not explain the between-method deficits on these matched loss cases.

\paragraph{Predicate design versus checker judgments.}
We review the saved task, predicate, and skill scope separately from the required-skill annotation: relevance does not establish predicate truth. Table~\ref{tab:app-routing-gap-attribution} restricts semantic attribution to the retrieved-but-rejected ranking-loss cases. Restrictive or misdirected predicates dominate this TheoremQA subset, whereas likely incorrect Boolean judgments dominate MedCalc-Bench. These are fallible retrospective assessments, not recovered checker explanations or experimentally established causes; ambiguous cases are retained explicitly.

\begin{table*}[t]
\centering
\small
\begin{tabular}{lrr}
\toprule
Assessment & MedCalc-Bench & TheoremQA \\
\midrule
Restrictive or misdirected predicate & 3 (8.57\%) & 39 (70.91\%) \\
Likely incorrect Boolean judgment & 25 (71.43\%) & 13 (23.64\%) \\
Ambiguous interpretation & 7 (20.00\%) & 3 (5.45\%) \\
\midrule
Retrieved-but-rejected ranking-loss cases & 35 & 55 \\
\bottomrule
\end{tabular}
\caption{Retrospective semantic assessment of required-skill rejections within the ranking-loss subset. Percentages use 35 MedCalc-Bench and 55 TheoremQA task--retriever cases, not all selection errors. The 39 restrictive TheoremQA cases constitute 59.09\% of its 66 total ranking-loss cases.}
\label{tab:app-routing-gap-attribution}
\end{table*}

For example, a TheoremQA question requests a P/E ratio from stock price and earnings per share, but the learned predicates additionally require dividend-discount-model inputs such as growth rates and payout ratios. The required skill ranks 22nd, sixth, and third under RouterBSI with BM25, Vector, and SkillSight, respectively, versus first under SkillRouter in each case. Conversely, a MedCalc-Bench question explicitly requests serum osmolality and supplies sodium, blood urea nitrogen, and glucose. Its predicate describes precisely this calculation from these inputs but evaluates false, leaving the required skill at rank 11 under BM25 versus first under SkillRouter. The latter supports a likely checker error rather than an overly narrow predicate.

\paragraph{Implications for refinement.}
Additional training or targeted human edits could address restrictive predicates, but neither intervention was tested here, and erroneous judgments of already appropriate predicates also require attention. Test results for the retained policies show mixed late progress: from iteration 5 to 10, MedCalc-Bench nDCG@10 changes by +1.713, 0, and 0 percentage points for BM25, Vector, and SkillSight, while TheoremQA changes by $-0.502$, +0.219, and $-0.542$ points. Thus, simply extending the same training loop is not an established remedy. Predicate edits and checker improvements should be evaluated separately using training/validation data before fresh held-out assessment.
